\chapter*{Nomenclature\markboth{}{Nomenclature}}
\addcontentsline{toc}{chapter}{Nomenclature}

% \vspace*{-3mm}
\subsubsection*{Variables and functions}

\begingroup
\renewcommand{\arraystretch}{1.2}
\renewcommand{\tabularxcolumn}[1]{p{#1}}
\begin{tabularx}{\textwidth}{@{}p{2.5cm}L}
    $p(x)$ & Probability density function with respect to variable $x$.\\
    $P(A)$ & Probability of event $A$ occurring.\\
    $\varepsilon$ & The Bayes error. \\
    $\varepsilon_u$ & The Bhattacharyya bound. \\
    $B$ & The Bhattacharyya distance. \\
    $s$ & An HMM state.  A subscript is used to refer to a particular state, e.g.\ $s_i$ refers to the $i^{\text{th}}$ state of an HMM. \\
    $\mathbf{S}$ & A set of HMM states. \\
    $\mathbf{F}$ & A set of frames. \\
    $\mathbf{o}_f$ & Observation (feature) vector associated with frame $f$. \\
    $\gamma_s(\mathbf{o}_f)$ & A posteriori probability of the observation vector $\mathbf{o}_f$ being generated by HMM state $s$. \\
    $\mu$ & Statistical mean vector. \\
    $\Sigma$ & Statistical covariance matrix. \\
    $L(\mathbf{S})$ & Log likelihood of the set of HMM states $\mathbf{S}$ generating the training set observation vectors assigned to the states in that set. \\
    $\mathcal{N}(\mathbf{x} | \mu, \Sigma)$ & Multivariate Gaussian PDF with mean $\mu$ and covariance matrix $\Sigma$.\\
    $a_{ij}$ & The probability of a transition from HMM state $s_i$ to state $s_j$. \\
    $N$ & Total number of frames or number of tokens, depending on the context. \\
    $D$ & Number of deletion errors. \\
    $I$ & Number of insertion errors. \\
    $S$ & Number of substitution errors. \\
\end{tabularx}
\endgroup


\newpage
\subsubsection*{Acronyms and abbreviations}

\begingroup
\renewcommand{\arraystretch}{1.2}
\begin{tabular}{@{}p{2.5cm} l}
    FTC     & Fault-Tolerant Control \\
    ADCS    & Attitude Determination and Control System \\
    LEO    & Low Earth Orbit \\
    COTS   & Commercial Off-The-Shelf \\
\end{tabular}
\endgroup

\subsection{Notation}
\begin{itemize}
    \item Vectors are denoted by bold lowercase letters, e.g.\ $\vec{x}$.
    \item Matrices are denoted by bold uppercase letters, e.g.\ $\vec{A}$.
    \item Scalars are denoted by lowercase letters, e.g.\ $a$.
\end{itemize}
It is also worth defining the following notation for sub and superscripts of vectors and matrices:
\begin{center}
% \centering
\begin{itemize}
    \item subscripts denote points, e.g.\ $\vec{r}_{AB}$ is the position vector from point $B$ to point $A$ ($A$ relative to $B$).
    \item superscripts denote frames, e.g.\ $\vec{v}^I_B$ is the velocity vector of point $B$ relative to inertial frame $I$ and $\vec{\omega}^{BI}$ is the angular velocity of frame $B$ relative to $I$.
    \item coordinates are denoted by superscript outside of brackets e.g.\ $[\vec{r}_{AB}]^I$ represents the position vector of point $A$ relative to $B$ expressed in coordinates of frame $I$.
    \item the $*^\times$ denotes the skew-symmetric matrix of a vector, e.g.\ $\vec{\omega}^\times = \begin{bmatrix} 0 & -\omega_z & \omega_y \\ \omega_z & 0 & -\omega_x \\ -\omega_y & \omega_x & 0 \end{bmatrix}$.
    \item the $\bar{*}$ denotes a unit vector, e.g.\ $\bar{\vec{v}} = \frac{\vec{v}}{||\vec{v}||}$.
    \item the $\hat{*}$ denotes the estimated value of a quantity, e.g.\ $\hat{\vec{\omega}}$ is the estimated angular velocity.
    \item the $\tilde{*}$ denotes the error of a quantity relative to an estimate, e.g.\ $\tilde{\vec{\omega}} = \vec{\omega} - \hat{\vec{\omega}}$ is the angular velocity estimation error.
\end{itemize}
\end{center}
These notations may be omitted when context makes one aspect or another clear. Other lowercase and upper case sub and superscripts may be used to denote other operators, quantities etc. which should be clear given context. 